\section{Conclusion}
\label{sec:conclusion}

We presented five mathematically grounded evaluation frameworks for large language models, drawing on persistent homology, sheaf theory, information geometry, topological data analysis, and random matrix theory.
Through a comprehensive survey of 11 evaluation dimensions and proof-of-concept implementations on \mmlu data, we demonstrated that these mathematical tools address fundamental gaps in current methodology: stability guarantees via the persistence stability theorem, composability analysis via sheaf cohomology, proper metric structure via Fisher-Rao geometry, structured failure detection via topological phase transitions, and contamination detection via spectral anomalies.

Our key finding is that mathematically rigorous evaluation is both feasible and informative.
The Wasserstein distance of 1.77 between model generations, consistency ratios ranging from 0.48 to 1.83, \fr distances up to 2.82, topological phase transitions at the 0.8 difficulty threshold, and Tracy-Widom statistics of 83.1 all provide quantitative insights inaccessible to standard accuracy metrics.
These are not merely alternative numbers---they carry formal guarantees that standard metrics fundamentally cannot provide.

\para{Future work.}
Three directions are immediate priorities.
First, applying these frameworks to real model outputs via API evaluation would validate the approach on ground-truth data rather than published scores.
Second, scaling topological analysis to 100+ models using approximate persistence algorithms would enable detection of richer topological features ($H_1$, $H_2$).
Third, developing an open-source toolkit integrating all five frameworks would enable the evaluation community to adopt mathematically grounded methods without requiring expertise in algebraic topology or spectral theory.
More broadly, we hope this work catalyzes a shift toward treating LLM evaluation as a formal mathematical discipline, with provable guarantees matching the rigor that the field's importance demands.
